% PV-Finder presentation — editable source for the complete deck.
%
% EDIT THIS FILE. All slide text, tables, diagrams and settings are here.
% Build: bash build.sh       Output: revised_deck.pdf
% Direct build: pdflatex slides.tex (run twice; produces slides.pdf).
%
% Figures are external PDF/PNG assets in figures/, research/ and ../v4/.
% To change text inside a generated plot, edit its Python script in review/.
% The previous split sources are preserved under review/multifile_source_*.
%
% SLIDE INDEX (search for "SLIDE NN" to jump to a frame)
%   01  Title
%   02  Motivation
%   03  One set of settings
%   04  HL-LHC: vertex separation at integral cuts 0.3 and 0.5
%   05  Unweighted vertex categories: four examples
%   06  Vertex categories: our unweighted definitions
%   07  PV-Finder and AMVF: vertex finding
%   08  PV-Finder + GNN: track-association diagnostics
%   09  Proposed paper structure
%   10  Planned plots and published HS benchmarks
%   11  Run 3 timing: the published comparison
%   12  HL-LHC: vertex counts versus pile-up
%   13  HL-LHC: track association versus pile-up
%   14  HL-LHC: vertex categories after track association
%   15  Target widths without AMVF: a simulation study
%   16  CPU vs GPU Comparison
%   17  Appendix
%   18  Backup: CPU and GPU timing breakdown
%   19  Timing Breakdown
%   20  Backup: Inference Timing: What Has Been Measured

% ============================================================================
% SETUP — fonts, theme, metadata and reusable formatting
% ============================================================================

\documentclass[xcolor=table,aspectratio=169]{beamer}

\usepackage[defaultsans,scale=0.95]{opensans}
\usepackage{FiraMono}
\usepackage{fourier}
\usefonttheme{professionalfonts}
\DeclareMathAlphabet{\mathcal}{OMS}{cmsy}{m}{n}
\newcommand\bmmax{2}
\usepackage{bm}
\usepackage[T1]{fontenc}

\usepackage{graphicx}
\usepackage{caption}
\captionsetup[figure]{labelformat=empty}
\usepackage{subcaption}
\usepackage{multicol}

\usepackage{tikz}
\usetikzlibrary{decorations.pathreplacing}
\usepackage{wrapfig}

% COLORS
\usepackage{xcolor}
\colorlet{myred}{red!80!black}
\colorlet{myblue}{blue!80!black}
\colorlet{mygreen}{green!60!black}
\colorlet{myorange}{orange!70!red!60!black}
\colorlet{mydarkred}{red!30!black}
\colorlet{mydarkblue}{blue!40!black}
\colorlet{mydarkgreen}{green!30!black}
\colorlet{mylightblue}{blue!60!white}
\colorlet{mydarkyellow}{yellow!70!black}
\colorlet{myyellowred}{yellow!70!red}

\usepackage{float}
\usepackage{hyperref}

\usepackage{amsmath, amsthm, amssymb, amsfonts}
\usepackage{esint}
\usepackage{mathtools}
\usepackage{listings}
\lstset{basicstyle=\ttfamily}

\usepackage{physics}
\usepackage[makeroom]{cancel}
% siunitx + mhchem commented out due to version clash with chemgreek on this TL2017 install
\usepackage[italic]{heppennames}
\usepackage{xpatch}
\makeatletter
\xpatchcmd\@HepConStyle
{\edef\@upcode{\updefault}}
{\ifdefined\shapedefault\edef\@upcode{\shapedefault}\else\edef\@upcode{\updefault}\fi}
{}{}
\makeatother
\usepackage{tensor}
\usepackage{feynmp-auto}
\usepackage{listofitems}
\usetikzlibrary{arrows.meta}
\usepackage[outline]{contour}
\usepackage{booktabs}

%%% ---- shared TikZ vocabulary for the Part I method figures ----
\usetikzlibrary{positioning,calc,backgrounds}
\tikzset{
    pvflow/.style   ={-{Latex[length=2.4mm,width=1.9mm]}, line width=0.9pt, draw=black!55},
    pvpoint/.style  ={-{Latex[length=2.0mm,width=1.6mm]}, line width=0.9pt},
    pvaxis/.style   ={-{Latex[length=1.9mm,width=1.5mm]}, line width=0.7pt, draw=black!60},
    pvdim/.style    ={{Latex[length=1.7mm,width=1.4mm]}-{Latex[length=1.7mm,width=1.4mm]},
                      line width=0.7pt},
    pvcurve/.style  ={line width=1.0pt, line join=round, line cap=round},
    pvthresh/.style ={dash pattern=on 2.6pt off 2.2pt, line width=0.8pt,
                      draw=myorange!85!black},
    pvsplit/.style  ={dash pattern=on 3pt off 2.4pt, line width=1.0pt, draw=pennred},
    pvguide/.style  ={line width=0.5pt, draw=black!15},
    pvcard/.style   ={rounded corners=6pt, draw=black!15, fill=black!2, line width=0.7pt},
    pvbox/.style    ={rounded corners=6pt, align=center, line width=0.9pt},
}
\newcommand{\pvcap}[1]{{\scriptsize\textcolor{black!55}{#1}}}


\addtobeamertemplate{frametitle}{}{\vspace{-15.5pt}}

%%%%%%%%%%%%%%%%%%%%%%%%%%%%%%%%%%%%%%

\newcommand{\backupbegin}{
    \newcounter{finalframe}
    \setcounter{finalframe}{\value{framenumber}}
    \setbeamertemplate{footline}{}
    \begin{frame}[plain]
        \centering
        \vfill
        {\color{pennred!28}\rule{0.30\linewidth}{1.1pt}}\\[0.8em]
        {\LARGE\textbf{\textcolor{pennred}{Backup}}}\\[0.8em]
        {\color{pennred!28}\rule{0.30\linewidth}{1.1pt}}
        \vfill
    \end{frame}
}
\newcommand{\backupend}{
    \setcounter{framenumber}{\value{finalframe}}
    \setbeamertemplate{footline}[frame number]
}

\newcommand{\ntrk}{\ensuremath{n_\mathrm{trk}}}

\def\thisframelogos{}
\newcommand{\framelogo}[1]{\def\thisframelogos{#1}}

\addtobeamertemplate{frametitle}{}{%
\begin{tikzpicture}[remember picture,overlay]
\node[anchor=north east,yshift=3.5pt] at (current page.north east) {%
    \foreach \img in \thisframelogos {%
        \hspace{.5ex}%
        \includegraphics[height=0.8cm, keepaspectratio]{\img}%
    }%
};
\end{tikzpicture}}

% Logo and historical-image paths are resolved through graphicspath below.

% Presentation mode is the default. Set talkbuild to 0 to display the extra
% reading notes already included in the source.
\providecommand{\talkbuild}{1}
\newcommand{\readingonly}[1]{\ifnum\talkbuild=0\relax#1\fi}
\newcommand{\talkonly}[1]{\ifnum\talkbuild=1\relax#1\fi}
\newcommand{\bnum}[2]{\ifnum\talkbuild=0 #1\else #2\fi}

% Part-divider helper: \partdivider{numeral}{title}{one-line subtitle}
\newcommand{\partdivider}[3]{%
    \begin{frame}[plain]
        \centering
        \vfill
        \begin{tikzpicture}
            \fill[pennred] (0,0) circle (0.72cm);
            \node[text=white, font=\Large\bfseries] at (0,0.02) {#1};
        \end{tikzpicture}\\[0.85em]
        {\LARGE\textbf{\textcolor{pennred}{#2}}}\\[0.75em]
        {\color{pennred!28}\rule{0.30\linewidth}{1.1pt}}\\[0.75em]
        {\small\textcolor{black!50}{#3}}
        \vfill
    \end{frame}%
}

%%%%%%%%%%%%%%%%%%%%%%%%%%%%%%%%%%%%%%

\usetheme{Philadelphia}
\hypersetup{colorlinks,allcolors=pennred}

\title[PV-Finder]{PV-Finder: End-to-end Primary Vertex Reconstruction}
\subtitle{Method, validation and results at high pile-up}

\author[Mattia Mauro]{Qi-Bin Lei, Rocky B. Garg, Lauren Tompkins,\\\textcolor{pennred}{Mattia Mauro}, Trey Warner}

\institute[]{Stanford University}

\date{September 25, 2026}

\titlegraphic{
    \hspace{5pt}
    \href{https://stanford.edu/group/stanford_atlas/}{\includegraphics[height=1cm]{logo-Stanford/Block_S_2_color.png}}
    \href{https://atlas.cern/}{\includegraphics[trim={0 11.5cm 0 11.5cm},clip,height=1cm]{logo-ATLAS/ATLAS-Logo-Ref-Pantone.pdf}}
    \href{https://iris-hep.org/}{\includegraphics[height=1cm]{logo-iris-hep/Iris-hep-4-no-long-name.png}}
}


% Small reusable text styles. Edit slide text in the numbered frames below.
\newcommand{\readingnote}[1]{\readingonly{\par\vspace{.35em}%
  {\footnotesize\color{black!65}#1\par}}}
\newcommand{\takeaway}[1]{\par\vspace{.65em}{\centering
  \setlength{\fboxsep}{7pt}\colorbox{pennred!5}{\parbox{.94\linewidth}{%
  \centering\small\textcolor{pennred}{#1}}}\par}}
\newcommand{\context}[1]{{\centering\footnotesize\color{black!55}#1\par}\vspace{.65em}}
\newcommand{\bigvalue}[2]{{\LARGE\bfseries\textcolor{pennred}{#1}}\\[.4em]{\small #2}\par}
\newcommand{\pubnote}{\href{https://cds.cern.ch/record/2858348}{ATL-PHYS-PUB-2023-011}}
\newcommand{\notelink}{\href{https://gitlab.cern.ch/atlas-physics-office/IDTR/ANA-IDTR-2026-01/ANA-IDTR-2026-01-INT1}{ANA-IDTR-2026-01-INT1}}
\newcommand{\treysource}{\href{https://gitlab.cern.ch/atlas-physics-office/IDTR/ANA-IDTR-2026-01/ANA-IDTR-2026-01-INT1/-/blob/4c50ed51500c897dd6a1b5bb52b0f1e3c12e7575/sections/04b_timing.tex}{Note Sec. 5.1 (18 Sep 2026)}}

\newcommand{\svtx}{\sigma_{\mathrm{vtx-vtx}}}
\setbeamertemplate{navigation symbols}{}
\hypersetup{pdftitle={PV-Finder: Method and Results, September 25, 2026},
  pdfauthor={Mattia Mauro, Qi-Bin Lei, Rocky B. Garg, Lauren Tompkins, Trey Warner},
  pdfsubject={Paper proposal, performance, width calibration and timing; revised September 24, 2026}}



\graphicspath{{../v4/}}

% Source links and small captions used on several slides.
\newcommand{\sourcecaption}[1]{\par\vspace{.5em}{\scriptsize\color{black!55}#1\par}}
\newcommand{\itkpaper}{\href{https://arxiv.org/html/2412.15090v2\#S5.SS5}{ITk performance paper, Sec.\ 5.5 (local draft pp.\ 27--29)}}
\newcommand{\runthreepaper}{\href{https://arxiv.org/abs/2308.09471}{Run 3 tracking paper, Fig.\ 13}}
\newcommand{\amvfnote}{\href{https://cds.cern.ch/record/2670380}{AMVF note (2019), Sec.\ 3.2, pp.\ 3--4}}

\begin{document}

\framelogo{}

% ============================================================================
% SLIDE 01 — Title
% ============================================================================
\begin{frame}[plain]
    \titlepage
\end{frame}

% ============================================================================
% SLIDE 02 — Motivation
% ============================================================================
\begin{frame}{Motivation}
    \relax
    \begin{columns}[c]
        \begin{column}{.53\linewidth}
            \textbf{Reconstruct primary vertices from tracks.}
            \par\vspace{1em}
            \begin{itemize}\setlength\itemsep{.7em}
                \item Nearby interactions overlap as pile-up increases.
                \item AMVF finds and fits vertices iteratively.
                \item PV-Finder predicts a vertex-density map; a GNN assigns the tracks.
            \end{itemize}
        \end{column}
        \begin{column}{.44\linewidth}\centering
            \bigvalue{$\mu\sim200$}{HL-LHC design pile-up}
            \par\vspace{1.5em}
            \begin{tikzpicture}[x=.70cm,y=.85cm]
                \draw[pvaxis] (0,0)--(7.3,0) node[right] {$z$};
                \foreach \x/\h in {.8/1.4,2.4/1.7,2.85/1.2,4.5/1.8,5/1.4,6.2/1.1}{
                \draw[pennred,line width=1.2pt] (\x,0)--(\x,\h);
                \fill[pennred] (\x,\h) circle (2pt);}
                \node[font=\scriptsize,text=black!55] at (3.5,-.65) {several vertices in the same event
                    (schematic)};
            \end{tikzpicture}
        \end{column}
    \end{columns}
    \readingnote{The current PU200 production uses full ITk acceptance. Trey Warner's Athena study measures
        the complete finder and GNN chain on the same events as AMVF, with CPU and GPU configurations. Earlier
        method development is described in \pubnote.}
\end{frame}

% Numerical provenance and qualifications: SOURCES.md and review/.


% ============================================================================
% SLIDE 03 — One set of settings
% ============================================================================
\begin{frame}{One set of settings}
    \vspace{.8em}
    \centering
    {\Large Integral threshold: \textbf{\textcolor{pennred}{0.3 everywhere}}}\par
    \vspace{1.8em}
    {\Large Vertex finding: \textbf{\textcolor{pennred}{$|\Delta z|\leq0.1$ mm everywhere}}}\par
    \vspace{1.8em}
    {\Large Vertex categories: \textbf{\textcolor{pennred}{close to ATLAS}}}\par
    \vspace{.45em}
    {\small Clean / Matched, Merged, Split, Fake}\par
    \vfill
    \sourcecaption{The ITk paper
        uses 0.1 mm for HS finding and we extend this to our eval. Current categories use unit track weights and an unlabelled-track proxy
        for Fake.}
\end{frame}

% ============================================================================
% SLIDE 04 — HL-LHC: vertex separation at integral cuts 0.3 and 0.5
% ============================================================================
\begin{frame}{HL-LHC: vertex separation at integral cuts 0.3 and 0.5}
    \centering
    \includegraphics[width=.99\linewidth,height=.77\textheight,keepaspectratio]{figures/resolution_integral_030_050.pdf}
\end{frame}

% Four independent examples of the agreed unweighted ntuple diagnostic.
% Every dot is one assigned track. The Split example includes both vertices;
% Fake is the unlabelled-track proxy, shown as a strict plurality below 50%.

% ============================================================================
% SLIDE 05 — Unweighted vertex categories: four examples
% ============================================================================
\begin{frame}{Unweighted vertex categories: four examples}
    \context{Each dot is one assigned track. Four independent examples of the convention used in our plots.}
    \centering
    \resizebox{.98\linewidth}{!}{%
    \begin{tikzpicture}[x=1cm,y=1cm,font=\small]
        \colorlet{catA}{mydarkblue!85}
        \colorlet{catB}{myorange!85!black}
        \colorlet{catC}{teal!70!black}
        \colorlet{catF}{black!45}
        \foreach \x/\w in {1.6/3.0,5.0/3.0,9.35/5.0,13.7/3.0}{
        \node[pvcard,minimum width=\w cm,minimum height=3.65cm] at (\x,2.05) {};
        }
        % Clean: one truth interaction contributes 5/6 of the assigned tracks.
        \node[font=\small\bfseries,text=mygreen!65!black] at (1.6,3.5) {Clean / Matched};
        \foreach \p in {(1.0,2.75),(1.6,2.75),(2.2,2.75),(1.0,2.12),(1.6,2.12)}
        \fill[catA] \p circle (5pt);
        \fill[catB] (2.2,2.12) circle (5pt);
        \node[align=center,font=\footnotesize] at (1.6,1.1) {$5/6\simeq83\%$ from one\\truth collision};
        % Merged: 60/40 avoids an ambiguous tie for the dominant truth origin.
        \node[font=\small\bfseries,text=myorange!85!black] at (5.0,3.5) {Merged};
        \foreach \p in {(4.4,2.75),(5.0,2.75),(5.6,2.75)}
        \fill[catA] \p circle (5pt);
        \foreach \p in {(4.7,2.12),(5.3,2.12)}
        \fill[catB] \p circle (5pt);
        \node[align=center,font=\footnotesize] at (5.0,1.1) {$60\%+40\%$\\Neither reaches $70\%$};
        % Split: both vertices are dominated by the same truth interaction.
        \node[font=\small\bfseries,text=myblue] at (9.35,3.5) {Split};
        \foreach \x/\name in {8.15/v_a,10.55/v_b}{
        \node[draw=black!25,rounded corners=3pt,minimum width=1.7cm,minimum height=1.1cm] at (\x,2.55) {};
        \node[font=\scriptsize,text=black!60] at (\x,3.08) {$\name$};
        \foreach \dx in {-.42,0,.42} \fill[catC] (\x+\dx,2.55) circle (5pt);
        }
        \node[font=\footnotesize] at (9.35,1.65) {$\sum p_T^2(v_a)>\sum p_T^2(v_b)$};
        \node[font=\footnotesize,text=mygreen!65!black] at (8.15,.98) {Keep Clean};
        \node[font=\footnotesize\bfseries,text=myblue] at (10.55,.98) {Relabel Split};
        % Fake proxy: 3 unlabelled + 2 from each truth collision.
        \node[font=\small\bfseries,text=myred] at (13.7,3.5) {Fake (proxy)};
        \foreach \p in {(13.1,2.75),(13.7,2.75),(14.3,2.75)}
        \fill[catF] \p circle (5pt);
        \foreach \p in {(13.05,2.12),(13.48,2.12)} \fill[catA] \p circle (5pt);
        \foreach \p in {(13.92,2.12),(14.35,2.12)} \fill[catB] \p circle (5pt);
        \node[align=center,font=\footnotesize] at (13.7,1.1) {$3/7$ unlabelled $>2/7$\\from either collision};
        % Colours refer to truth origin within each independent example.
        \fill[catA] (1.0,-.45) circle (5pt);
        \node[anchor=west,font=\footnotesize] at (1.25,-.45) {truth collision A};
        \fill[catB] (4.7,-.45) circle (5pt);
        \node[anchor=west,font=\footnotesize] at (4.95,-.45) {truth collision B};
        \fill[catC] (8.4,-.45) circle (5pt);
        \node[anchor=west,font=\footnotesize] at (8.65,-.45) {same truth collision};
        \fill[catF] (12.65,-.45) circle (5pt);
        \node[anchor=west,font=\footnotesize] at (12.9,-.45) {unlabelled tracks};
    \end{tikzpicture}%
    }
    \par\vspace{.65em}\small
    Track counts set the category; $\sum p_T^2$ only resolves the Split label.
    \sourcecaption{Unweighted approximation to \amvfnote. The paper uses fit weights and truth-classified fake
        tracks; neither is available here. Main results: top-1 GNN assignment, reco
        $\ntrk\geq2$. ATLAS labels a reconstructed vertex fake when the summed fit weight of fake tracks exceeds the contribution from every individual true collision.}
\end{frame}

% ============================================================================
% SLIDE 06 — Vertex categories: our unweighted definitions
% ============================================================================
\begin{frame}{Vertex categories: our unweighted definitions}
    \context{Every assigned track counts once. Truth labels identify its collision of origin.}
    \small\renewcommand{\arraystretch}{1.5}
    \begin{tabular}{p{3.2cm}p{10.3cm}}
        \toprule
        \textbf{Clean / Matched} & At least 70\% of the assigned tracks come from one true collision.\\
        \textbf{Merged} & A true collision is the largest contribution, but none reaches 70\% of the assigned
            tracks.\\
        \textbf{Split} & Several reco vertices have the same dominant true collision. Only the vertex with the
            largest $\sum p_T^2$ keeps Clean or Merged.\\
        \textbf{Fake (proxy)} & The unlabelled-track group is the largest contribution.\\
        \bottomrule
    \end{tabular}
    \par\vspace{.6em}\footnotesize
    Both methods retain reco vertices with at
        least two stored tracks. Unlabelled tracks
        enter the category fractions.
    \sourcecaption{Adapted from \amvfnote; Track and Vertex Reconstruction with the ATLAS Inner Detector,
        Sec.\ 6.5, pp.\ 44--45 (local draft). Those papers use fit weights and truth-classified fake tracks.}
\end{frame}

% ============================================================================
% SLIDE 07 — PV-Finder and AMVF: vertex finding
% ============================================================================
\begin{frame}{PV-Finder and AMVF: vertex finding}
    \context{$14$ TeV $t\bar t$ simulation, $\langle\mu\rangle=192.5$, $|\eta|<4$; same 1,920 events; integral
        cut $0.3$.}
    \centering
    \small\renewcommand{\arraystretch}{1.12}
    \begin{tabular}{lccc}
        \toprule
        Measurement & $z$ window & PV-Finder & AMVF\\
        \midrule
        All-vertex positional efficiency & $\pm0.1$ mm & $74.11\pm0.10\%$ & $70.99\pm0.10\%$\\
        Efficiency after GNN + selection & $\pm0.1$ mm & $71.07\pm0.10\%$ & $70.99\pm0.10\%$\\
        Unmatched reco / event, after selection & $\pm0.1$ mm & $15.05\pm0.09$ & $25.46\pm0.11$\\
        \midrule
        Hard-scatter finding efficiency & $\pm0.1$ mm & $99.90^{+0.07}_{-0.14}\%$ &
            $99.79^{+0.10}_{-0.16}\%$\\
        \bottomrule
    \end{tabular}
    \par\vspace{.25em}\raggedright\footnotesize
    \textbf{Efficiency:} truth PVs with $\ntrk\geq2$ selected reco tracks and a reco within $0.1$ mm / all
        such truth PVs. One reco may recover several truth PVs.\\[.2em]
    \textbf{Track association:} top-1 GNN target per track; keep reco $\ntrk\geq2$, no reassignment.\\[.2em]
    \textbf{Unmatched/event:} retained reco vertices with no eligible truth PV within $0.1$ mm, divided by
        the number of events. This is a count per event, not a percentage.\\[.2em]
    \textbf{Hard scatter:} events with any vertex within $0.1$ mm of the true HS position / all events.
    \sourcecaption{The first and HS rows use PV-Finder peaks before GNN; the middle two rows apply the full-pileup GNN
        (epoch 213) and reco selection. AMVF uses reco $\ntrk\geq2$ throughout. HS criterion: ITk p.\ 27.}
\end{frame}

% ============================================================================
% SLIDE 08 — PV-Finder + GNN: track-association diagnostics
% ============================================================================
% Fit the longer diagnostic notes on one page without removing text.
\begin{frame}[shrink=5]{PV-Finder + GNN: track-association diagnostics}
    \context{Same 1,920 events, $185\leq\mu\leq200$; integral cut $0.3$; reco $\ntrk\geq2$
        for both.}
    \centering\small\renewcommand{\arraystretch}{1.2}
    \begin{tabular}{lcc}
        \toprule
        Measurement & PV-Finder + GNN & AMVF (stored links)\\
        \midrule
        Track efficiency & $69.06\pm0.06\%$ & $68.32\pm0.06\%$\\
        Assignment purity & $69.69\pm0.06\%$ & $68.65\pm0.06\%$\\
        Assigned fraction & $99.10\pm0.01\%$ & $99.52\pm0.01\%$\\
        \midrule
        Clean / reconstructed vertices & $49.75\pm0.13\%$ & $35.17\pm0.10\%$\\
        \bottomrule
    \end{tabular}
    \par\vspace{.25em}\raggedright\footnotesize
    \textbf{Truth-labelled tracks:} efficiency = correct / all input tracks; purity = correct / assigned
        links; assigned fraction = assigned / all input tracks.\\[.3em]
    Correct: the target's dominant unweighted truth origin equals the track's origin. All truth origins enter,
        including one-track origins; efficiency counts each track once.\\[.3em]
    Clean uses the unit-weight approximation, including Split relabelling.
    \takeaway{Full-pileup GNN: $+0.75\pm0.02$\ efficiency points and $+1.05\pm0.02$\ purity points over AMVF.}
\end{frame}

% ============================================================================
% SLIDE 09 — Proposed paper structure
% ============================================================================
\begin{frame}{Proposed paper structure}
    \begin{columns}[T]
        \begin{column}{.45\linewidth}
            \small
            \textbf{1\quad Introduction}\\[.45em]
            Why vertex reconstruction matters at high pile-up.\\[1em]

            \textbf{2\quad ATLAS detector}\\[.45em]
            2.1\quad Run 3 Inner Detector\\[.3em]
            2.2\quad The ITk upgrade\\[1em]

            \textbf{3\quad Data and MC samples}\\[.45em]
            Run 3 data and simulation; HL-LHC simulation.\\[.35em]
            Track selection, pile-up and beam spot.
        \end{column}

        \begin{column}{.51\linewidth}
            \small
            \textbf{4\quad PV-Finder and the track-to-vertex GNN}\\[.45em]
            Inputs, both architectures, training and chosen settings.\\[1em]

            \textbf{5\quad Performance}\\[.45em]
            5.1\quad Run 3 simulation\\[.3em]
            5.2\quad Run 3 data\\[.3em]
            5.3\quad HL-LHC simulation\\[.3em]
            5.4\quad Timing studies\\[1em]

            \textbf{6\quad Conclusion}\\[.45em]
            Main improvements and remaining limitations.
        \end{column}
    \end{columns}

    \takeaway{Describe both algorithms in Section 4; compare their performance with AMVF in Section 5.}
    \sourcecaption{Section order follows the professor's outline, dated 23 September 2026.}
\end{frame}

% ============================================================================
% SLIDE 10 — Plots in the current paper draft
% ============================================================================
\begin{frame}{Plots in the current paper draft}
    \small

    \textbf{Performance figures already specified in the paper skeleton}

    \begin{itemize}
        \item \textbf{Fig.\ 3a:} Number of reconstructed vertices:
        PV-Finder vs.\ AMVF in simulated $t\bar t$ events,
        classified by vertex type.

        \item \textbf{Fig.\ 3b:} Number of reconstructed PV-Finder vertices
        versus $\langle\mu\rangle$, classified by vertex type.

        \item \textbf{Fig.\ 4:} Longitudinal separation between reconstructed
        PV-Finder and AMVF primary vertices in simulated $t\bar t$ events.

        \item \textbf{Figs.\ 5--6:} Additional performance plots currently
        carrying placeholder copies of the Fig.\ 3--4 captions.

        \item \textbf{Fig.\ 7:} Processing time per event versus
        $\langle\mu\rangle$ for PV-Finder and AMVF.
    \end{itemize}

    \vspace{.4em}
    \textbf{Also in the algorithm section}
    \begin{itemize}
        \item Fig.\ 1: PV-Finder algorithm diagram.
        \item Fig.\ 2: Track-to-vertex GNN algorithm diagram.
    \end{itemize}

    \takeaway{The immediate task is to populate the figures already defined
    in the paper skeleton and clarify the intended content of Figs.\ 5--6.}

    \sourcecaption{PV-Finder paper draft, 23 September 2026, Figs.\ 1--7.}
\end{frame}


% ============================================================================
% SLIDE 12 — HL-LHC: vertex counts versus pile-up
% ============================================================================
\begin{frame}{HL-LHC: vertex counts versus pile-up}
    \centering
    \includegraphics[width=.99\linewidth,height=.86\textheight,keepaspectratio]{figures/finder_and_categories_vs_mu.pdf}
\end{frame}

% ============================================================================
% SLIDE 13 — HL-LHC: track association versus pile-up
% ============================================================================
\begin{frame}{HL-LHC: track association versus pile-up}
    \centering
    \includegraphics[width=.98\linewidth,height=.83\textheight,keepaspectratio]{figures/ttva_efficiency_purity_vs_mu.pdf}
\end{frame}

% ============================================================================
% SLIDE 14 — HL-LHC: vertex categories after track association
% ============================================================================
\begin{frame}{HL-LHC: vertex categories after track association}
    \centering
    \includegraphics[width=.96\linewidth,height=.82\textheight,keepaspectratio]{figures/ttva_categories_no_cut_vs_mu.pdf}
\end{frame}

% ============================================================================
% SLIDE 15 — Target widths without AMVF: a simulation study
% ============================================================================
\begin{frame}{Target widths without AMVF: a simulation study}
    \context{Truth supplies the track groups; a separate fit estimates the vertex position.}
    \small
    \textbf{Group the reconstructed tracks by their truth vertex, with $\ntrk\geq2$.}\\[.3em]
    Fit reconstructed measurements. Use the truth position only to measure the residual afterwards.
    \par\vspace{.2em}
    \textbf{Start with a simple one-dimensional fit.}\\[.3em]
    At a fixed transverse beam position, a weighted least-squares estimate is
    \[
    \hat z=\frac{\sum_i z_{0,i}/\sigma_{z_0,i}^2}{\sum_i1/\sigma_{z_0,i}^2}.
    \]
    All $z_0$ values must refer to the same beam line. Check this approximation and the residual bias. A full
        3D fit needs a standalone vertex fitter and full track covariances.
    \par\vspace{.2em}
    \textbf{Measure the width of $\hat z-z_{\rm truth}$ versus track multiplicity.}\\[.3em]
    Fit a positive curve to set the Gaussian training-target width. Validate it on separate events.
    \takeaway{A fit is still needed. Truth grouping replaces AMVF finding and track assignment.}
    \sourcecaption{Proposal only. Truth grouping omits association mistakes. Target width is not the final
        position error.}
\end{frame}

% Restored from V4, with its original numerical results.

% ============================================================================
% SLIDE 16 — CPU vs GPU Comparison
% ============================================================================
\begin{frame}{CPU vs GPU Comparison}
    \context{Same events, from input tracks to a vertex container with associated tracks.}
    \begin{columns}[c]
        \begin{column}{.68\linewidth}\centering
            \includegraphics[width=\linewidth,height=\bnum{.57}{.67}\textheight,keepaspectratio]{images/athena_cpu_gpu_amvf.png}
        \end{column}
        \begin{column}{.29\linewidth}\small
            \textbf{One lxplus-gpu node}\\[.5em]
            Intel Xeon Silver 4216\\
            @ 2.10 GHz\\[.5em]
            NVIDIA Tesla T4\\[1em]
            \textbf{Near PU200}\\[.5em]
            \textcolor{pennred}{\Large$\sim9\times$ faster}\\[.4em]
            PV-Finder + GNN on GPU versus AMVF on CPU.
        \end{column}
    \end{columns}
    \par\vspace{.45em}{\scriptsize\color{black!55}V4 benchmark: elapsed time around Athena execution; 200--500
        events/sample. Settings not remeasured here. \treysource.}
    \readingnote{The note describes four Run 3 collision-data and six Run 4 simulation samples; the original
        plot labels the collision-data series Run 2/3. The figure is reproduced unchanged. The ninefold
        comparison is reported in the note; the PV-Finder CPU implementation is slower than AMVF in this
        benchmark.}
\end{frame}

\appendix

% ============================================================================
% SLIDE 17 — Appendix
% ============================================================================
\partdivider{A}{Appendix}{Timing measurements from V4}
% Restored V4 backup measurements; no new timing claim.

% ============================================================================
% SLIDE 18 — Backup: CPU and GPU timing breakdown
% ============================================================================
\begin{frame}{Backup: CPU and GPU timing breakdown}
    \context{597 events across three Run 4 PU200 samples; mean wall time per complete event.}
    \begin{columns}[c]
        \begin{column}{.59\linewidth}\centering\footnotesize
            \begin{tabular}{lrr}\toprule
                Stage & CPU (ms) & GPU (ms)\\\midrule
                Track selection & 30.54 & 29.96\\
                Tensor building & 0.10 & 0.10\\
                Finder inference & 2188.67 & 37.56\\
                Peak finding & 0.05 & 0.05\\
                Vertex building & 0.29 & 0.22\\
                GNN track association & 181.65 & 8.30\\
                \quad Graph building & 2.83 & 2.50\\
                \quad GNN inference & 177.89 & 5.20\\
                \quad Edge selection & 0.68 & 0.41\\\midrule
                \textbf{Total} & \textbf{2401.28} & \textbf{76.18}\\\bottomrule
            \end{tabular}
        \end{column}
        \begin{column}{.37\linewidth}\small
            \bigvalue{$31.5\times$}{faster PV-Finder + GNN with CUDA}
            \par\vspace{.9em}\raggedright
            \textbf{GPU time:}\\[.4em]
            Finder + GNN inference: $56.1\%$.\\[.5em]
            Track selection: $39.3\%$.\\[.7em]
            Input preparation becomes a major part of the total.
        \end{column}
    \end{columns}
    \par\vspace{.8em}{\scriptsize\color{black!55}Same Xeon Silver 4216 / Tesla T4 setup. GNN substages sum to
        8.11 ms; reported total: 8.30 ms. Table from \treysource.}
    \readingnote{The GPU column is the complete CPU+CUDA execution path, including CPU track selection and
        graph construction. GNN association comprises graph building (2.50 ms), inference (5.20 ms), and edge
        selection (0.41 ms); components are rounded. The 31.5 ratio uses 2401.28/76.18 and compares the two
        PV-Finder implementations, not AMVF.}
\end{frame}

% ============================================================================
% SLIDE 19 — Timing Breakdown
% ============================================================================
\begin{frame}{Timing Breakdown}
    \centering
    \includegraphics[width=.99\linewidth,height=.82\textheight,keepaspectratio]{image.png}
\end{frame}

% ============================================================================
% SLIDE 20 — Backup: Inference Timing: What Has Been Measured
% ============================================================================
\begin{frame}{Backup: Inference Timing: What Has Been Measured}
    \relax
    \begin{columns}[c]
        \begin{column}{.59\linewidth}\includegraphics[width=\linewidth,height=\bnum{0.488}{.61}\textheight,keepaspectratio]{images/timing.png}
        \end{column}
        \begin{column}{.38\linewidth}
            \centering\bigvalue{4.5 ms}{per 40 mm $z$ interval, earlier model}
            \par\vspace{1.2em}\raggedright\small
            NVIDIA A100, batch size one.\\[.6em]
            Inference only; $0\leq\mu\leq80$.\\[.6em]
            Mean time changes little across this range.
        \end{column}
    \end{columns}
    \takeaway{Earlier model: inference only, one of twelve intervals; not a full-event
        timing.}\readingnote{Each sub-event is a 40 mm interval along the beam axis. The study timed only the
        first of twelve intervals per event: 12,150 calls, 150 per integer pile-up bin. Input loading, GPU
        transfer and peak finding are outside the timer. The 1.2 ms spread is the standard deviation across
        calls; plotted errors are the standard error on each bin mean. Separate profiler operator sums are not
        a wall-clock breakdown or a measured batching gain.}
    \par\vspace{.35em}{\scriptsize\color{black!55}Input loading, GPU transfer and peak finding are outside the
        timer. Original V4 measurement.}
\end{frame}

\end{document}
